\begin{claim}{D.2.32}{Proved}{The coordinated once-suspended map stops at the eleventh skeleton}
Use the spaces and the two-local cellular models of D.2.30, and put
$R_i=K_i/K_6$. There is a specified ordinary map
\[
 g_9:R_9\longrightarrow W=\Omega\fib(\Sigma^2K(\FF_2,2)
                         \to K(\FF_2,4))
\]
with the degree-eight values of D.2.30, whose composite from $K_9$
pulls $\vartheta_7$ back to $\alpha_3\operatorname{Sq}^1\alpha_3$.
It extends over $R_{10}$, but no such extension extends over $R_{11}$.
For the normalized eleven-cell relation described below, every
ten-cell filling has the same nonzero obstruction
\[
 p\in\pi_{10}W\cong(\ZZ/2)^4,
 \qquad h_{\ZZ}(p)=4N',\qquad \Psi(p)=0.
 \tag{D.2.32a}
\]
Here $h_{\ZZ}$ is computed in the equivalent group $\pi_9E$ of
D.2.22, $N'$ generates its order-eight homology summand, and
$\Psi$ is the suspended cone detector of D.2.23.

The lower map is fixed explicitly in the proof. This does not
classify other lower maps even when they vanish on $K_6$
(for example, a different choice on the $J$-Moore top), nor
maps nonzero on $K_6$. It does not prove that the coefficient
in D.2.28 always vanishes or assert survival through later
nullification stages.
\end{claim}
\begin{proof}
All undisplayed homology coefficients are $\FF_2$; integral
calculations are two-local. Retain the source labels of
D.2.19--D.2.20 and abbreviate $A_0,B_0,C_5,D_6,E_6,F_7,X_7$
by $A,B,C,D,E,F,X$ within this proof. Set $R=Q^4A$,
$P=AC$, $N=Q^5A$, and $J=AB+R$ in mod-two degree seven.

\medskip\noindent
\textbf{The ordinary range and the fixed lower map.}
The source $R_i$ is six-connected. Its maps to $W$ lift to the
five-connected cover $W'$; write $E=\Omega W'$ as in D.2.22.
The evaluation $\Sigma E\to W'$ is eleven-connected: its
homotopy fiber is $E*E$, which is ten-connected. Thus relative
lifts through dimension eleven are available. Moreover $\Sigma E$
is five-connected, so its tenth homotopy group is stable. All
source attachments used below are in the ordinary stable range
of a six-connected source. These facts justify the finite stable
cell calculations with their specified ordinary maps; no
unrestricted stable replacement of the mapping problem is used.

Collapse $K_6$ in the established source eight- and nine-cell
tables. Before the nine-cells, the factors are
\[
 M(\ZZ/2,7)_F\vee M(\ZZ/4,7)_J\vee S^7_X\vee S^7_R
           \vee S^8_H\vee S^8_P\vee S^8_N.
\]
The nine-attachments, in the previous unit notation $I_9,J_9,L_9$,
are
\[
\begin{array}{c|c@{\qquad}c|c}
 I_9&\eta F&J_9&2H+\eta F\\
 L_9&2H&AD&\eta J+\eta R\\
 AE&0&BC&2P+\eta J+\eta R\\
 Q^6A&2N&Q^5B&\eta R\\
 A^3&0&&
\end{array}
 \tag{D.2.32b}
\]
The former even-$\nu$ and lower torsion ambiguities disappear
on collapsing $K_6$. The remaining integral degrees and second
squares detect every component in this range. Cell slides give
two factors on which the chosen map will be nonzero:
\[
 \begin{aligned}
 J_*&=M(\ZZ/4,7)_J\cup_{i\eta}e^9_U,
                 &&U=AD-Q^5B,\\
 P_*&=M(\ZZ/2,8)_P,
                 &&\text{top cell }BC-AD.
 \end{aligned}
 \tag{D.2.32c}
\]
The other factors are
$S^7_R\cup_\eta e^9_{Q^5B}$,
$M(\ZZ/2,7)_F\cup_{i\eta}e^9_{I_9}$,
$M(\ZZ/2,8)_H$, $M(\ZZ/2,8)_N$, $S^7_X$, and the
three free nine-spheres with labels $J_9-I_9-L_9,AE,A^3$.

Use the target $L_8$ of D.2.22. Send $J_*$ to $\Sigma L_8$
by the Moore extension of $\rho$ and the nullhomotopy of
$\rho\eta$ supplied by its $c$-cell. Send $P_*$ to the
suspended target $R$-Moore factor. Send every other factor
to zero. Composing with evaluation defines the specified
$g_9$. Its Moore top has zero Hurewicz, its $P$ value is
$R_8$, and its $N$ value is zero. These are exactly the
coordinated choices of D.2.30, including the original
$B^2$ value $R_8$. In particular it extends over dimension
ten by that theorem. The specified nine-map already lands
in $\Sigma L_8$, a fact needed to control integral fillings.

\medskip\noindent
\textbf{All integral degree-ten cycles.}
Let $m,n,o,t$ denote the four unit ten-classes dual to
$u^3v,u^2w,v^3,t_9$. The complete degree-ten basis is
\[
 AR,Q^6B,Q^7A,A^2B,AF,AX,BD,BE,C^2,m,n,o,t.
\]
Its nonzero first Bocksteins are
\[
\begin{gathered}
 \beta(AR)=\beta(A^2B)=A^3,\qquad
 \beta(Q^6B)=Q^5B,\qquad \beta(AX)=AD+AE,\\
 \beta(BD)=AD,\qquad\beta(BE)=AE,\qquad
 \beta(o)=J_9+L_9.
\end{gathered}
 \tag{D.2.32d}
\]
They have rank five. Define seven cycles
\[
 \begin{array}{lll}
 K_1=AR+A^2B,&K_2=C^2,&K_3=Q^7A,\\
 K_4=AF,&K_5=m+n,&K_6=t,\\
 K_7=AX+BD+BE.&&
 \end{array}
 \tag{D.2.32e}
\]
The full degree-eleven Bockstein table is equivalently the
following list of its eight nonzero entries, together with
the nine listed zero entries:
\[
\begin{array}{c|cccccccc}
 z&BR&Q^6C&Q^8A&AG&BF&\sigma(u^2v^2)^\vee
                         &\sigma(w^2)^\vee&BX\\ \hline
 \beta z&K_1&K_2&K_3&K_4&K_4&K_5&K_6&K_7
\end{array}
 \tag{D.2.32f}
\]
The zero entries are
$AQ^5A,Q^7B,A^2C,AH,AB^2,CD,CE,
\sigma(u^5)^\vee,\sigma(uvw)^\vee$.
Thus no nonprimitive degree-eleven source class is omitted.
Cartan proves the product entries, and
$\beta Q^s=(s-1)Q^{s-1}$ proves the Dyer entries, including
$Q^5C=C^2$. The unit entries follow from Serre's polynomial
algebra and $\operatorname{Sq}^1t_9=w^2$
\cite[\S2, nos.~8--9]{I_Serre1953}
\cite[Chapter~I, Theorem~1.1]{I_CohenLadaMay1976}.
These are operations inside the source, not operations
preserved by arbitrary ordinary maps.

The known $H_8A=(\ZZ/2)^3$ and the mod-two dimensions
$\dim H_9A=9$, $\dim H_{10}A=13$ show by UCT that
$H_9A$ has six cyclic summands and $H_{10}A$ has seven.
All these integral groups are finite two-primary groups.
Rank seven in (D.2.32f) forces every summand of $H_{10}A$
to have exponent two. Consequently
\[
 H_{10}(A;\ZZ_{(2)})\cong(\ZZ/2)^7.
 \tag{D.2.32g}
\]
The integral cycle lattice of $R_{10}$ is free of rank
seven, with a basis reducing to (D.2.32e). Its eleven-boundary
image is exactly twice the lattice. Odd-unit changes and
even cell additions therefore normalize seven eleven-cell
boundaries to $2K_1,\ldots,2K_7$ and every other boundary
to zero, while preserving all mod-two labels.

The last relation can be retained exactly. Choose the
elementary integral chains with
$dB=2A$, $dD=4C$, $dE=-4C$, $dX=2(D+E)$.
The integral nine-cycle labels used in (D.2.32b) are
$AD'=AD+2BC$ and $AE'=AE-2BC$: thus
$d(BD)=2AD'$, $d(BE)=2AE'$, and
$d(AX)=-2(AD'+AE')$. Their reductions are the displayed
mod-two labels $AD,AE$; the following product-chain
identity is unchanged by this normalization.
The actual Pontryagin chain product gives
\[
 d(AX+BD+BE)=0,\qquad d(BX)=2(AX+BD+BE)=2K_7.
 \tag{D.2.32h}
\]
A cellular chain comparison transports these identities
to the finite model. The nonzero mod-two class makes
$K_7$ a primitive two-local cycle coordinate. This
normalization does not assume that $K_7$ is spherical
before its own ten-cell is attached.

\medskip\noindent
\textbf{The Moore coextension detected by the seventh cycle.}
The complete second-square table on the seven cycles is
\[
 \operatorname{Sq}_*^2K_i=0\ (i=1,2,3,4,6),\qquad
 \operatorname{Sq}_*^2K_5=G_8+H_8,\qquad
 \operatorname{Sq}_*^2K_7=P+B^2.
 \tag{D.2.32i}
\]
In (D.2.32c), $P+B^2$ is precisely the $J$-Moore top:
the two $P$ coordinates cancel. All seven cycles have
zero integral ninth-cell degrees.
The Chang cofiber sequence gives
\[
 \pi_9J_*\cong\ZZ_{(2)}\{\gamma\}\oplus\ZZ/2\{\delta\},
 \qquad h(\gamma)=2U,
\]
where $\delta$ is the image of the Moore eta-coextension
and is detected by the second square from the Moore top.
This is also the explicit generator of
\cite[Proposition~4.1 and Table~3]{I_Li2024}.
Thus $K_7$ projects to exactly $\delta$, while the other
six cycles project to zero. Their projections to $P_*$
are all zero: $\pi_9P_*$ is the bottom eta line, and
none has the corresponding second-square coefficient.
All other source factors were assigned the zero map.
There is no cross Whitehead term; such terms begin in
dimension thirteen for these source factors.

Here is the exact target image of $\delta$, with all
indeterminacies. Desuspend once and use
\[
 D_0=\cofib\left(
 \begin{pmatrix}4&\eta\\0&2\end{pmatrix}:
 S^5\vee S^6\longrightarrow S^5\vee S^6\right)
\]
from D.2.22. Its bottom classes satisfy
$2\rho=-\alpha\eta$; write $n=\alpha\nu$, of order four.
Let $\mu:M(\ZZ/4,6)\to D_0$ extend $\rho$.
The Moore comparison $f:M(\ZZ/2,6)\to M(\ZZ/4,6)$ has
bottom degree two and top degree one
\cite[(3.1)--(3.2)]{I_Li2024}.
The maps $\mu f$ and $-\alpha\bar\eta$ agree on the bottom,
so their difference factors through the top seven-sphere.
The group $\pi_7D_0$ is generated by $b,c$, both of order two.
Compose with $\widetilde\eta_2:S^8\to M(\ZZ/2,6)$.
The actual extension/coextension computation of
\cite[\S2, p.~191]{I_Gray2014} gives the stable bracket
$\bar\eta\widetilde\eta_2\in\{2\nu,6\nu\}$.
Both choices have image $2n$. Therefore
\[
 \mu\widetilde\eta_4=2n
       \pmod{\langle b\eta,c\eta\rangle},
 \qquad \mu\widetilde\eta_4=2n\text{ in }\pi_8L_8.
 \tag{D.2.32j}
\]
Changing the Moore coextension by bottom $\eta^2$ adds
$\rho\eta^2=c\eta$, also killed in $L_8$.
Thus the displayed equality includes every relevant
extension and coextension ambiguity.

\medskip\noindent
\textbf{Every eleven-cell obstruction, including the cone detector.}
The first six cycle attachments are null already in
$\Sigma L_8$, so choose their fillings there. They contribute
no target degree-ten homology. The seventh has composite
$2n$ by (D.2.32j); fill it with two copies of the normalized
target $N'$-cell, which attaches by $n$ in $L_9$.
This reference filling has integral coordinate $2N'$.
The map and the homotopy identifying its attachment with
$2n$ both lie in $\Sigma L_8$, so no additional top
homology coordinate is introduced.

It follows from the complete boundary normalization that
every eleven-cell obstruction has zero integral Hurewicz,
except the $BX$ coordinate, whose Hurewicz is $4N'$.
The only integral-Hurewicz kernel is $\langle s\rangle$,
$s=\rho\nu$, by D.2.22. The actual suspended cone
detector removes this last ambiguity on every cell.

For the chosen lower map the pullbacks, on the original
source, are
\[
 q_7=\alpha\beta,\qquad \ell_9=\alpha\delta,
 \qquad \beta=\operatorname{Sq}^1\alpha,
        \delta=\operatorname{Sq}^2\beta.
\]
The extra $B^2$-dual in the $\mathcal R_8$ pullback has
zero first Bockstein and does not alter $\ell_9$.
With $\gamma=\operatorname{Sq}^2\alpha$ and
$X^*=\operatorname{Sq}^3\beta$, Cartan gives
\[
 \operatorname{Sq}^4q_7=
       \gamma\delta+\beta X^*
       =\operatorname{Sq}^2\ell_9.
 \tag{D.2.32k}
\]
In particular the cross term
$\operatorname{Sq}^1\alpha\operatorname{Sq}^1\delta$
in the second expression has been retained.
Although a degree-seven quotient class can have a relative
ambiguity, $H^{11}(R_{11})\to H^{11}(K_{11})$ is an
isomorphism. Thus (D.2.32k) also determines the top value
on each actual quotient-cell attachment.

For an attaching composite $\omega:S^{10}\to W$, the
classes $\vartheta_7,\lambda_9$ extend uniquely over its
actual cone. The top value of
$\operatorname{Sq}^4\vartheta_7+\operatorname{Sq}^2\lambda_9$
is zero on the normalized $p$ and one on $s$.
There is an actual map from the corresponding source
one-cell attachment to this cone with relative top degree
one. The top source coordinate is nonzero modulo two,
since its integral cellular boundary is even. Equation
(D.2.32k) makes the pulled-back detector zero. Hence
every obstruction has zero detector: all zero-Hurewicz
obstructions vanish, and the $BX$ obstruction is exactly
$p$, rather than $p+s$.
The normalization here is explicitly
$h_{\ZZ}(p)=4N'$ and $\Psi(p)=0$; these conditions remove
its possible $s$ ambiguity.

Finally, changing the ten-cell filling of $K_i$ by
$z_i\in\pi_{10}W$ changes its normalized relation by
$2z_i$, and changes every zero-boundary relation by zero.
This is the actual cell-coaction law; in this range no
Whitehead term involving the added ten-sphere can occur.
D.2.31 proves
\[
                  2\pi_{10}W=0.
 \tag{D.2.32l}
\]
Thus no choice of any ten-cell filling removes the
nonzero $BX$ obstruction. All other eleven-cell
obstructions have already been accounted for. This
proves the claimed nonextension for the specified
coordinated nine-map, with the stated limitations.
\end{proof}
